\documentclass[10pt,a4paper]{article}
\usepackage[utf8]{inputenc}
\usepackage[T1]{fontenc}
\usepackage[margin=18mm]{geometry}
\usepackage{longtable}
\usepackage{hyperref}
\usepackage{enumitem}
\setlist[itemize]{noitemsep, topsep=2pt}
\setlength{\parindent}{0pt}
\begin{document}

\section*{Non-Verbal \& Multimodal Discourse Analysis Worksheet}

\textit{Qualitative analysis of pauses, trailing sentences, prosody/intonation, vocalisations (coughs, sighs, scoffs), and kinesics (hand gestures, facial expressions, posture) in political news discourse}\par\smallskip

\bigskip\hrule\bigskip

\subsection*{Guidelines: Five non-verbal \& paralinguistic features to look for in every event}

\begin{longtable}{|p{0.23\textwidth}|p{0.23\textwidth}|p{0.23\textwidth}|p{0.23\textwidth}|}
\hline
Feature & What to look for & Illustrative example (from sample) & Questions to answer \\ \hline \endhead
\textbf{1. Pauses, Silences \& Abrupt Stops} & Micro-pauses \texttt{(.)}, timed silences \texttt{(1.5)}, pregnant/dramatic pauses before a loaded word, hesitant stumbles, abrupt self-interruptions or cut-offs (\texttt{---}), and forced silences (muting a panelist's mic or talking over them). & Anchor pauses \texttt{(1.5)} before saying \textit{"'loota'"} for dramatic weight; Addl. DCP cuts off mid-sentence (\textit{"wo toh--- abhi investigation chal raha hai"}) when asked if police struck first; panelist's mic faded during counter-question. & Where does the pause or stop occur? Is it deliberate/rhetorical (for suspense, emphasis, or gravitas) or defensive/evasive (cognitive load, self-repair, avoiding blame)? Who controls turn-taking and silence? \\ \hline
\textbf{2. Trailing Sentences \& Unfinished Utterances} & Clauses that trail off (\texttt{...}), drop in volume at the end, or leave the conclusion unspoken (aposiopesis) so the audience supplies the implication without the speaker taking explicit responsibility. & \textit{"Aur phir wahan jo log ghuse the... aap samajh hi sakte hain..."}; \textit{"500 salon ke sangharsh ke baad agar aisa..."} trailing off with a slow head shake. & What exact thought is left unfinished? What presupposition or shared bias is the viewer invited to fill in? Does trailing off act as a safe hedge to imply higher-level complicity or "anti-national" motive without stating it? \\ \hline
\textbf{3. Intonation, Pitch, Volume \& Prosody} & Pitch contour (rising \texttt{$\uparrow$} for incredulity/challenge; falling \texttt{$\downarrow$} for finality), volume shifts (shouting vs. hushed whisper), emphatic word stress, speech rate (rapid-fire urgency vs. slow solemnity), and tonal register (sarcastic, reverential, combative). & Hushed, reverential cadence on \textit{"Prabhu Shri Ram ke Mandir"} vs. sharp, high-pitched interrogation on \textit{"par Akhilesh ji Ram Mandir kyon$\uparrow$ nahin jate?!"}; satirical sing-song inflection on \textit{"priventiv arenjment"}. & Which exact words receive pitch peaks or volume stress? How does vocal tone turn neutral words into irony, mockery, or moral outrage? Does the speaker shift prosodic register depending on which actor is discussed? \\ \hline
\textbf{4. Coughs, Sighs \& Non-Lexical Vocalisations} & Audible sighs (\texttt{[sigh]}), scoffs/snorts (\texttt{[scoff]}), ironic chuckles (\texttt{[chuckle]}), throat-clearing/coughs (\texttt{[cough]}), sharp intakes of breath/gasps (\texttt{[gasp]}), and tongue-clicks (\texttt{[tsk]}). & Heavy \texttt{[sigh]} before reporting Trust oversight lapses (signalling reluctant sorrow rather than hostility); mocking \texttt{[scoff]} or \texttt{[chuckle]} when reading an opponent's denial; \texttt{[cough / throat-clear]} when pivoting away from an awkward fact. & What emotion or epistemic stance does the vocalisation perform (exasperation, moral disgust, disbelief, ridicule, discomfort)? Does a sigh or scoff pre-judge a speaker's claim before any verbal rebuttal is made? \\ \hline
\textbf{5. Hand Gestures, Facial Expressions \& Posture (Kinesics)} & Hand movements (air-quotes, finger-pointing/wagging, desk-thumping, open palms, counting on fingers, dismissive wave), facial cues (raised eyebrows, eye-rolls, smirk, furrowed brow), head nods/shakes, leaning into camera, and visual props (bandaged limbs, documents). & Anchor forms physical air-quotes on \textit{"'loota'"} or \textit{"'peaceful' protest"}; host points finger/pen at lens during an ultimatum; injured DCP holds bandaged hand up to camera; creator raises one eyebrow and leans in for a conspiratorial reveal. & How do gestures reinforce, contradict, or sharpen the spoken words? During split-screen listening, does the host's face (nodding vs. smirking/eye-rolling) cue the audience on whom to trust? How are bodily visuals used as evidence? \\ \hline
\end{longtable}

\subsection*{Quick Multimodal Transcription Notation Key (use inside transcript excerpts)}

\begin{longtable}{|p{0.15\textwidth}|p{0.15\textwidth}|p{0.15\textwidth}|p{0.15\textwidth}|p{0.15\textwidth}|p{0.15\textwidth}|}
\hline
Symbol & Meaning & Example in transcript & Symbol & Meaning & Example in transcript \\ \hline \endhead
\texttt{(.)} & Brief micro-pause (< 0.5 sec) & \texttt{woh (.) suparavijan lacking tha} & \texttt{$\uparrow$ / $\downarrow$} & Marked rising / falling pitch & \texttt{kyon$\uparrow$ FIR nahin karaee$\downarrow$} \\ \hline
\texttt{(1.5)} & Timed pause in seconds & \texttt{kisne (1.5) 'loota'?} & \texttt{WORD / \_word\_} & Loud volume / emphatic stress & \texttt{ye PAAP hai / \_bilkul\_ galat} \\ \hline
\texttt{---} & Abrupt stop, cut-off, or self-repair & \texttt{police ki taraf se--- abhi jaanch...} & \texttt{$^\circ$word$^\circ$} & Softened / whispered delivery & \texttt{$^\circ$ander ki baat ye hai$^\circ$} \\ \hline
\texttt{...} & Trailing off / unfinished sentence & \texttt{aap jaante hain ki kaun log...} & \texttt{[sigh] [cough] [scoff]} & Non-lexical vocalisations & \texttt{[sigh] 500 salon ke baad...} \\ \hline
\texttt{[overlap]} & Simultaneous / overlapping speech & \texttt{[Anjana: suniye---] [Guest: nahin!]} & \texttt{\{\{gesture: ...\}\}} & Hand gesture / facial cue / posture & \texttt{\{\{air-quotes on 'loota'\}\}} \\ \hline
\end{longtable}

\textbf{Methodological Reminder:} A non-verbal cue is not automatically manipulative or ideological---a cough, pause, or rapid cadence may reflect studio audio delay, live time pressure, or a speaker's habitual style. Always check (1) whether the cue clusters with verbal features, and (2) whether the speaker's delivery shifts noticeably across different actors or events (within-outlet baseline).\par\smallskip

\bigskip\hrule\bigskip

\subsection*{Event 1. Ram Mandir donation controversy (Non-Verbal \& Paralinguistic Layer)}

\begin{itemize}
  \item \textbf{Focus:} How do prosody, pauses, trailing sentences, vocalisations (sighs, coughs, scoffs), and body language frame the controversy---as solemn devotional sorrow, administrative scrutiny, or partisan combat?
  \item \textbf{Outlet:} Aaj Tak (\textit{Black \& White} / Special Bulletin with Anjana Om Kashyap) | \textbf{Date:} 23 July 2024
  \item \textbf{Video title:} Ram Mandir Donation Controversy / Counting Theft \& Audit Scrutiny Bulletin | \href{https://youtu.be/B4L1pz0o4dc}{https://youtu.be/B4L1pz0o4dc} | \textbf{Timestamp(s):} \texttt{[00:00]}--\texttt{[08:12]}
  \item \textbf{Analyst:} Baani Solanki (\texttt{2025114016}) | \textbf{Language:} Hindi (transcribed in Roman script)
  \item \textbf{Format / Visual setup:} \texttt{[ ]} TV Panel Debate \texttt{[x]} TV Bulletin \texttt{[ ]} Solo YouTube Explainer \texttt{[ ]} Interview / Discussion
  \item \textbf{Primary speaker(s) observed:} Anjana Om Kashyap (standing/seated studio news anchor delivering a dramatic prime-time monologue with CCTV footage inserts, temple visuals, and opposition soundbites)
\end{itemize}

\subsubsection*{Event-specific non-verbal questions (Event 1)}

\noindent\textbf{1.} \textbf{Vocal \& gestural marking of core labels ("scam", "loota", "paap", "gadabdi", "avyavastha"): When the controversy is named, what non-verbal cues accompany the label? (Look for physical air-quotes, vocal scare-quotes, dramatic pauses before the word, hushed/reverential drops in volume, or shouted stress.)}\par\smallskip
\begin{itemize}
  \item \textbf{Notes:}
Anjana Om Kashyap's vocal and gestural marking of core labels is the clearest example in our dataset of \textbf{sacralising individual theft while sanitising institutional failure}:\par\smallskip
  \item \textbf{Moral/Religious Labels (\texttt{"paap"}, \texttt{"aparadh"}):} When describing the counting clerk stealing cash from the offering boxes (\texttt{[00:42]}), she pauses \texttt{(1.2)}, drops her voice to a hushed, trembling register \texttt{$^\circ$...$^\circ$}, and then stresses \texttt{"PAAP"} with a deep, mournful falling pitch \texttt{$\downarrow$} and a slow head shake: \textit{"Prabhu Shri Ram ke charnon mein chadhaye gaye daan ki chori (1.2) ye sirf aparadh nahin (.) \_PAAP\_ hai$\downarrow$"} (\texttt{[00:48]}).
  \item \textbf{Political Accusation Labels (\texttt{"'loota'"}, \texttt{"'ghotala'"}):} When quoting Samajwadi Party leaders (Akhilesh Yadav / Awadhesh Prasad) who alleged that the Temple Trust \textit{"looted"} donations (\texttt{[04:15]}), Anjana pauses \texttt{(1.5)} before the word \textit{"'loota'"} and forms \textbf{physical air-quotes with both hands} \texttt{\{\{gesture: bilateral two-finger air-quotes on 'loota'\}\}} accompanied by raised eyebrows and a sharp, incredulous rising tone \texttt{$\uparrow$}, visually quarantine-marking the word \textit{"loota"} as a partisan slander.
  \item \textbf{Institutional Lapses (\texttt{"avyavastha"}, \texttt{"supervision lacking"}):} When describing the Trust's own lack of frisking and audit protocols (\texttt{[02:35]}), she switches to a brisk, low-arousal administrative cadence with a brief micro-pause \texttt{(.)}: \textit{"Wahan par (.) thoda supervision lacking tha, thodi \_avyavastha\_ thi$\downarrow$"} (\texttt{[02:42]}), accompanied by a small, minimising hand wave \texttt{\{\{gesture: small horizontal hand wave\}\}}.
\end{itemize}

\noindent\textbf{2.} \textbf{Differential demeanour toward actors (Temple, Trust leadership, counting clerks, donors, opposition accusers): How do the speaker's voice pitch, facial expression, and posture shift when referring to Prabhu Shri Ram / devotees vs. top figures (Champat Rai, Nripendra Misra, CM Yogi) vs. arrested counting staff vs. opposition leaders (Akhilesh Yadav, Awadhesh Prasad)?}\par\smallskip
\begin{itemize}
  \item \textbf{Notes:}
Anjana exhibits four sharply differentiated non-verbal personas depending on the actor named:\par\smallskip
  \item \textbf{Prabhu Shri Ram \& Devotees (\texttt{[00:15]}--\texttt{[01:10]}):} Hushed, reverential cadence (\texttt{$^\circ$Prabhu Shri Ram ke Mandir$^\circ$}), slightly bowed head, softened eyes, and hands brought together near her chest/waist in a semi-devotional posture \texttt{\{\{gesture: hands clasped gently at waist + bowed head\}\}}.
  \item \textbf{Arrested Counting Clerk / Bank Staff (\texttt{[01:20]}--\texttt{[02:20]}):} Stern, moralising scowl, pointing a rigid index finger toward the CCTV screen showing the clerk stuffing notes into his pocket (\texttt{[01:45]}), with clipped, disapproving downward stress.
  \item \textbf{Top Figures (Nripendra Misra, Champat Rai, CM Yogi, \texttt{[03:10]}--\texttt{[03:55]}):} Deferential, matter-of-fact institutional tone; upright posture and respectful nods when noting that Nripendra Misra ordered a digital counting overhaul.
  \item \textbf{Opposition Leaders (Akhilesh Yadav, Awadhesh Prasad, \texttt{[04:10]}--\texttt{[06:45]}):} Combative, high-pitched prosecutorial interrogation \texttt{$\uparrow$}, chin raised, right index finger wagging at the camera lens \texttt{\{\{gesture: assertive index-finger wag at lens\}\}} (\textit{"Par Akhilesh ji Ram Mandir kyon$\uparrow$ nahin jaate?!"}, \texttt{[05:18]}).
\end{itemize}

\noindent\textbf{3.} \textbf{Solemn devotional grief vs. high-decibel partisan spectacle: How do pauses, sighs, speech rate, and hand gestures construct the segment's emotional register? Where is delivery slow, hushed, and mournful (sighs, bowed head, folded hands), and where does it escalate into rapid-fire shouting, desk-thumping, or mocking irony?}\par\smallskip
\begin{itemize}
  \item \textbf{Notes:}
The bulletin is structured as a dramatic two-act non-verbal performance:\par\smallskip
  \item \textbf{Act I --- Solemn Devotional Grief (\texttt{[00:00]}--\texttt{[03:50]}):} Speech rate is slow (\textasciitilde{}115 WPM), punctuated by heavy audible \texttt{[sigh]} vocalisations (\texttt{[00:35]}, \texttt{[02:30]}), long \texttt{(1.2)}--\texttt{(1.5)} pauses, mournful head shakes, and hushed volume \texttt{$^\circ$...$^\circ$} lamenting that after \textit{"500 salon ke sangharsh ke baad..."} (\texttt{[00:38]}) a petty clerk betrayed devotees' faith in the sanctum.
  \item \textbf{Act II --- High-Decibel Partisan Counter-Attack (\texttt{[04:00]}--\texttt{[07:50]}):} The moment the Samajwadi Party's allegations are introduced, Anjana's speech rate jumps to \textasciitilde{}155 WPM, her pitch rises sharply \texttt{$\uparrow$}, her volume surges, and her hand gestures shift from mournful clasped hands to aggressive finger-pointing and palm-chopping against opposition "hypocrisy" (\texttt{"Jo log karsevakon par goli chalane walon ke saath the..."}, \texttt{[05:40]}).
\end{itemize}

\noindent\textbf{4.} \textbf{Non-verbal evaluation of financial scrutiny: When questions about unreceipted offerings, lack of frisking, or land deals are raised, how does the host/speaker react non-verbally (earnest nodding, leaning into the camera, open-palm appeals vs. eye-rolls, scoffs, dismissive hand waves, or incredulous rising pitch $\uparrow$)?}\par\smallskip
\begin{itemize}
  \item \textbf{Notes:}
Anjana splits financial scrutiny into two categories:\par\smallskip
  \item \textbf{Clerical Oversight (Lack of Frisking / Manual Counting):} Treated with \textbf{reluctant, maternal sorrow}---she lets out a heavy \texttt{[sigh]} at \texttt{[02:30]}, furrows her brow, and opens her palms (\texttt{"sawal toh banta hai ki wahan frisking kyon nahin hui$\downarrow$"}, \texttt{[02:48]}), framing the lapse as innocent naivety (\textit{"bholapan"} / \textit{"avyavastha"}) by a trusting religious body.
  \item \textbf{Systemic Corruption Allegations by Opposition:} Met with \textbf{incredulous rising pitch \texttt{$\uparrow$}, mocking \texttt{[scoff]} vocalisations, and dismissive hand waves} (\texttt{[04:25]}, \texttt{[05:10]}), treating broader audit demands by SP/Congress leaders as malicious attempts to defame Hindu faith.
\end{itemize}

\noindent\textbf{5.} \textbf{Performing certainty vs. hedging and innuendo (pauses, stops, trailing off): Are allegations delivered with definitive falling intonation $\downarrow$, finger-counting, and pointing at on-screen evidence (CCTV, documents), or hedged with hesitant pauses (.), abrupt self-interruptions (---), throat-clearing [cough], and trailing sentences (...) that imply higher-level complicity without naming it?}\par\smallskip
\begin{itemize}
  \item \textbf{Notes:}
  \item \textbf{Certainty on Individual Clerk's Guilt:} Delivered with definitive falling intonation \texttt{$\downarrow$} and direct pointing at the red-circled CCTV footage (\texttt{"CCTV mein range haathon pakda gaya$\downarrow$"}, \texttt{[01:52]}).
  \item \textbf{Hedging on Trust Leadership's Responsibility:} Delivered with micro-pauses \texttt{(.)}, throat-clearing \texttt{[cough]}, and euphemistic self-repair \texttt{---}: \textit{"Trust ki taraf se--- matlab jo counting agency aur bank ke log the (.) unki monitoring mein thodi kami rahi..."} (\texttt{[02:38]}), deflecting blame away from Champat Rai or Nripendra Misra onto external bank staff.
  \item \textbf{Innuendo against Opposition via Trailing Off (\texttt{...}):} At \texttt{[05:55]}, she trails off with a narrowed gaze: \textit{"Aur aaj wahi log Ram Mandir ke chande par sawaal utha rahe hain jinhone hamesha Mandir ka... aap jaante hain ki inka asli maksad kya hai..."} (\texttt{[06:02]}), leaving the accusation of "anti-Hindu conspiracy" unspoken for the viewer to fill in.
\end{itemize}

\noindent\textbf{6.} \textbf{Turn-taking, interruptions, and acoustic treatment of denials: In debates/interviews: who receives calm, uninterrupted speaking time and sympathetic nods, and who is cut off mid-sentence (---), talked over, or muted? In solo videos: how are clips of official denials framed via reaction pauses, raised eyebrows, or sarcastic chuckles?}\par\smallskip
\begin{itemize}
  \item \textbf{Notes:}
In this studio bulletin, recorded soundbites are framed asymmetrically:\par\smallskip
  \item \textbf{Official Trust / Government explanations} (Nripendra Misra's assurance of automated counting machines, \texttt{[03:25]}) are introduced with a calm, reassuring nod and followed by a supportive summary (\texttt{"ab poori vyavastha high-tech ki ja rahi hai$\downarrow$"}, \texttt{[03:45]}).
  \item \textbf{Opposition soundbites} (Awadhesh Prasad / SP leaders demanding a high-level inquiry, \texttt{[04:30]}) are cut short after 8 seconds, immediately followed by Anjana's mocking \texttt{[scoff]}, raised eyebrows, and a rapid-fire counter-interrogation directed at the screen (\texttt{"Sawaal poochhne se pehle Akhilesh Yadav ye batayein---"}, \texttt{[04:42]}).
\end{itemize}

\noindent\textbf{7.} \textbf{Performing the boundary between religious devotion and institutional critique: How does the presenter use vocal tone, sighs, and body language to signal personal reverence for the Temple (or secular objectivity) while criticising---or defending---the Trust's financial management?}\par\smallskip
\begin{itemize}
  \item \textbf{Notes:}
Anjana explicitly performs \textbf{devotional insiderhood}: she never critiques the Temple management from a secular distance, but rather speaks \textit{as an anguished devotee}. Before mentioning the counting lapse at \texttt{[02:28]}, she lets out an audible \texttt{[sigh]}, lowers her voice \texttt{$^\circ$...$^\circ$}, and places her hand over her heart \texttt{\{\{gesture: right hand pressed to heart + bowed head\}\}}: \textit{"Yeh khabar batate hue bhi (.) mann ko dukh hota hai..."} (\texttt{[02:32]}). This non-verbal performance signals that her mild critique of the counting arrangement comes from love for Lord Ram, insulating the channel from any charge of being anti-Mandir while simultaneously delegitimising outsider/Opposition criticism as hostile.\par\smallskip
\end{itemize}

\noindent\textbf{8.} \textbf{Self-presentation as the "guardian of public trust": How does the speaker use direct eye-gaze, posture, moral sighs, or emphatic gestures (pointing at the viewer, hands on heart/desk) to position themselves as the authentic protector of devotees or public accountability?}\par\smallskip
\begin{itemize}
  \item \textbf{Notes:}
Anjana positions herself as the \textbf{moral guardian of the pious Hindu devotee (\texttt{shraddhalu})} against both petty thieves inside the counting room and political opportunists outside it. She uses intense, unblinking eye-gaze into the teleprompter, places her hand on her chest when invoking \textit{"karodon Ram bhakton ki aastha"} (\texttt{[01:05]}), and points her pen/index finger sharply at the lens (\texttt{[05:18]}, \texttt{[07:10]}) when challenging opposition politicians, visually casting herself as the shield of the Temple's sanctity.\par\smallskip
\end{itemize}

\subsubsection*{Apply the five non-verbal features (Event 1)}

\begin{longtable}{|p{0.31\textwidth}|p{0.31\textwidth}|p{0.31\textwidth}|}
\hline
Non-verbal feature & Timestamp + spoken anchor words + non-verbal notation & What it does in context (stance, framing, implication) \\ \hline \endhead
\textbf{1. Pauses, silences \& stops} & \texttt{[02:38]} \textit{"Trust ki taraf se--- matlab jo bank ke log the (.) unki monitoring mein kami rahi"} / \texttt{[04:15]} \textit{"SP ne kaha ki chanda (1.5) 'loota' gaya$\uparrow$"} & The abrupt self-repair \texttt{---} at \texttt{[02:38]} shifts blame from the Trust to bank staff; the \texttt{(1.5)} pause before \textit{"'loota'"} isolates the opposition's word for ridicule. \\ \hline
\textbf{2. Trailing sentences (\texttt{...})} & \texttt{[00:38]} \textit{"[sigh] \texttt{500} salon ke sangharsh ke baad agar Mandir mein aisa..."} / \texttt{[06:02]} \textit{"Jinhone hamesha Mandir ka... aap samajh hi sakte hain..."} & At \texttt{[00:38]}, trailing off with a head shake evokes unspeakable devotional sorrow; at \texttt{[06:02]}, it insinuates that opposition critics harbour an anti-Ram agenda without stating it outright. \\ \hline
\textbf{3. Intonation, pitch \& volume} & \texttt{[00:48]} \textit{"$^\circ$Prabhu Shri Ram ke charnon mein$^\circ$... ye \_PAAP\_ hai$\downarrow$"} vs. \texttt{[05:18]} \textit{"Par Akhilesh ji Ram Mandir kyon$\uparrow$ nahin jaate?!"} & Hushed reverential cadence \texttt{$^\circ$...$^\circ$} on the Temple contrasts sharply with high-pitched, combative rising interrogation \texttt{$\uparrow$} aimed at Akhilesh Yadav. \\ \hline
\textbf{4. Coughs, sighs \& vocalisations} & \texttt{[02:30]} \texttt{[[sigh]]} \textit{"Mann ko dukh hota hai ki wahan suparavijan lacking tha..."} / \texttt{[04:40]} \texttt{[[scoff]]} (right after SP leader's byte) & The heavy \texttt{[sigh]} frames Trust oversight failures as reluctant family sorrow; the \texttt{[scoff]} instantly invalidates the Opposition's call for an audit. \\ \hline
\textbf{5. Hand gestures, face \& posture} & \texttt{[00:25]} \texttt{\{\{gesture: hands clasped at waist + bowed head\}\}} / \texttt{[04:15]} \texttt{\{\{gesture: bilateral physical air-quotes on 'loota'\}\}} / \texttt{[05:18]} \texttt{\{\{gesture: index-finger wag at lens\}\}} & Clasped hands perform piety; physical air-quotes mark \textit{"loota"} as fake news; finger-wagging turns the bulletin into an indictment of the Opposition. \\ \hline
\end{longtable}

\subsubsection*{Short summary (Event 1 Non-Verbal Synthesis)}

\begin{itemize}
  \item \textbf{What overall interpretation or emotional stance does the non-verbal delivery encourage?}
The non-verbal delivery encourages \textbf{reverential sorrow over an isolated clerical sin (\texttt{"paap"}) paired with fierce partisan indignation against opposition critics}. Viewers are guided to view the theft as the work of one greedy bank clerk, the Trust's lapse as innocent administrative oversight (\textit{"avyavastha"}), and the Opposition's questions as malicious political opportunism.\par\smallskip
  \item \textbf{Whose voice or physical presence is acoustically/visually foregrounded, undermined, or silenced?}
\textbf{Anjana's devotional-prosecutorial persona and the Trust's corrective measures} are foregrounded with calm authority and reverential pauses. \textbf{The arrested clerk} is visually shamed via CCTV loops, while \textbf{Opposition leaders (SP/Congress)} are acoustically and visually undermined through physical air-quotes \texttt{\{\{air-quotes on 'loota'\}\}}, mocking \texttt{[scoff]} vocalisations, and rapid counter-questions.\par\smallskip
  \item \textbf{What evidence supports your non-verbal reading? What complicates it (e.g., speaker baseline, format)?}
The unmistakable acoustic shift between Act I (\textasciitilde{}115 WPM, hushed \texttt{$^\circ$...$^\circ$}, \texttt{[sigh]}, clasped hands on Temple/Trust) and Act II (\textasciitilde{}155 WPM, loud \texttt{WORD}, rising pitch \texttt{$\uparrow$}, finger-wagging on Akhilesh Yadav) proves a strong within-speaker shift across actors. While high theatricality is part of Hindi prime-time TV bulletin conventions, the directional asymmetry of air-quotes and self-repairs (\texttt{Trust ki taraf se--- matlab bank ke log}, \texttt{[02:38]}) clearly protects institutional leadership.\par\smallskip
\end{itemize}

\bigskip\hrule\bigskip

\subsection*{Event 2. Students' protest regarding NEET paper leak (Non-Verbal \& Paralinguistic Layer)}

\begin{itemize}
  \item \textbf{Focus:} How do prosody, pauses, trailing sentences, vocalisations, and visual body language construct the legitimacy of student protesters, police use of force, and state accountability?
  \item \textbf{Outlet:} Aaj Tak (Live Prime-Time Debate / Ground \& Studio Broadcast with Anjana Om Kashyap) | \textbf{Date:} 2024 (NEET Protest \& Parliament March Live Coverage)
  \item \textbf{Video title:} NEET Paper Leak Protest Live: Students' Parliament March, Police Clash \& Political Debate | \href{https://www.youtube.com/live/JXLIs0WtJAM}{https://www.youtube.com/live/JXLIs0WtJAM} | \textbf{Timestamp(s):} \texttt{[00:00]}--\texttt{[32:40]}
  \item \textbf{Analyst:} Baani Solanki (\texttt{2025114016}) | \textbf{Language:} Hindi (transcribed in Roman script)
  \item \textbf{Format / Visual setup:} \texttt{[x]} TV Panel Debate \texttt{[ ]} TV Bulletin \texttt{[ ]} Solo YouTube Explainer \texttt{[x]} Interview / Discussion
  \item \textbf{Primary speaker(s) observed:} Anjana Om Kashyap (studio anchor), Additional DCP / Police Officer (live guest showing bandaged hand), Student / Youth Wing Protesters, and BJP vs. Opposition Spokespersons in multi-window split-screen
\end{itemize}

\subsubsection*{Event-specific non-verbal questions (Event 2)}

\noindent\textbf{1.} \textbf{Prosodic \& gestural marking of participant labels ("students/bachche" vs. "protesters" vs. "upadravi/goons"): How do vocal tone, volume, facial expression, and hand gestures shift across labels? (e.g., softened, empathetic pitch for "students" vs. harsh, staccato stress, scowls, or accusatory finger-pointing on "upadraviyon", "goons", or "anti-social elements".)}\par\smallskip
\begin{itemize}
  \item \textbf{Notes:}
Anjana performs a sharp \textbf{acoustic and gestural bifurcation} between "genuine students" at home and "protesters/rowdies" on the street:\par\smallskip
  \item When referring to \textit{"asli chhatra"} or \textit{"masoom bachche"} studying at home (\texttt{[03:15]}), her voice softens slightly with a concerned nod.
  \item However, the moment she refers to those marching toward Parliament and clashing at barricades (\texttt{[05:40]}, \texttt{[11:20]}, \texttt{[18:50]}), she shifts to the labels \textit{"pradarshankari"}, \textit{"chhatra sangathan ke karyakarta"}, and \textit{"upadravi"} (rowdies), delivering \textit{"UPADRAVIYON"} with \textbf{harsh, staccato stress, a deep scowl, and an accusatory pen-point toward the riot B-roll} \texttt{\{\{gesture: sharp pen-point at screen + furrowed scowl\}\}} (\texttt{"Kya chhatron ke bhes mein \_upadraviyon\_ ne kanoon haath mein liya$\uparrow$?!"}, \texttt{[06:12]}).
\end{itemize}

\noindent\textbf{2.} \textbf{Delivery of students' stated demands: When the protesters' core demands (Education Minister's resignation, compensation, release of Sonam Wangchuk, exam reform) are mentioned, are they delivered slowly with deliberate pauses and finger-counting, rushed through in a flat monotone, mocked with air-quotes/scoffs, or trailed off (...)?}\par\smallskip
\begin{itemize}
  \item \textbf{Notes:}
When protesters' broader political demands (such as the Education Minister's resignation or release of detained activists/Sonam Wangchuk alongside NEET re-examination, \texttt{[14:10]}--\texttt{[15:30]}) are mentioned, Anjana \textbf{rushes through them in a rapid, dismissive list} before stopping on the political demands with \textbf{incredulous rising pitch \texttt{$\uparrow$}, raised eyebrows, and a mocking \texttt{[scoff]}}: \textit{"Ab NEET ke naam par istifa maang rahe hain (.) Sonam Wangchuk ki baat kar rahe hain... matlab mudda exam hai ya \_siyasat\_$\uparrow$?!"} (\texttt{[14:45]}), accompanied by a two-handed dismissive wave \texttt{\{\{gesture: bilateral dismissive hand wave\}\}}.\par\smallskip
\end{itemize}

\noindent\textbf{3.} \textbf{Where does vocal \& kinetic energy peak---grievance, street disruption, or electoral horse-race?: Track where the speaker's volume, pitch range, pace, and gesture intensity reach their climax: during accounts of the paper leak/student suicides, during street clashes and barricade breaches, or during partisan BJP-vs-Opposition sparring.}\par\smallskip
\begin{itemize}
  \item \textbf{Notes:}
Vocal and kinetic energy peaks overwhelmingly during \textbf{street clashes/barricade breaches (\texttt{[05:30]}--\texttt{[12:40]}) and high-decibel BJP-vs-Opposition studio sparring (\texttt{[19:00]}--\texttt{[30:00]})}:\par\smallskip
  \item While initial references to the NEET paper leak are delivered at a moderate pace, Anjana's volume, pitch, and gesture intensity skyrocket when showing barricades being pushed and police officers injured (\texttt{"Dekhiye kaise \_pathar\_ chalaaye gaye! Barricade \_toda\_ gaya!"}, \texttt{[07:15]}), leaning forward over her desk and striking her palm with her pen \texttt{\{\{gesture: striking palm with pen\}\}}.
\end{itemize}

\noindent\textbf{4.} \textbf{Non-verbal attribution of responsibility (sighs, stops, head shakes, pointing): When examination authorities (NTA), the Education Minister, police, opposition politicians, or students are named as responsible, what non-verbal cues accompany the attribution (exasperated sighs, mocking laughs, abrupt self-stops ---, or direct finger-pointing at the lens)?}\par\smallskip
\begin{itemize}
  \item \textbf{Notes:}
  \item \textbf{Students / Opposition Youth Wings:} Named as responsible for street chaos with \textbf{direct finger-pointing at their split-screen windows}, stern head shakes, and loud interruptions (\texttt{"Aapke karyakartaon ne police par haath uthaya$\downarrow$!"}, \texttt{[12:10]}).
  \item \textbf{Police / Administration:} Completely absolved non-verbally via \textbf{sympathetic head nods and pained facial expressions} (\texttt{[08:20]}--\texttt{[10:45]}).
  \item \textbf{NTA / Government:} When an opposition panelist tries to attribute primary blame to the Education Minister (\texttt{[21:15]}), Anjana lets out an exasperated \texttt{[sigh]} and cuts them off mid-sentence \texttt{---} (\texttt{"CBI jaanch toh ho rahi hai na--- phir sadak par hungama kyon$\uparrow$?!"}, \texttt{[21:22]}).
\end{itemize}

\noindent\textbf{5.} \textbf{Prosody of questions---open inquiry vs. loaded rhetorical challenge: Are questions about the protest delivered with level, open intonation, or with sharp rising pitch $\uparrow$, sarcastic stress, raised eyebrows, and trailing clauses (...) that presuppose protesters are violent or politically manipulated?}\par\smallskip
\begin{itemize}
  \item \textbf{Notes:}
Questions directed at student leaders and opposition spokespersons are \textbf{100\% loaded rhetorical challenges} delivered with sharp rising pitch \texttt{$\uparrow$}, sarcastic stress, arched eyebrows, and accusatory trailing clauses (\texttt{...}), presupposing guilt:\par\smallskip
  \item To student/opposition representative (\texttt{[11:45]}): \textit{"Shantipurna pradarshan mein police walon ke haath kaise$\uparrow$ toot gaye?! Kya ye chhatra hain ya phir inke peechhe..."} (\texttt{[11:52]}), trailing off with a pointed stare that presupposes political conspiracy.
  \item Conversely, questions to the Additional DCP (\texttt{[08:30]}) are delivered with a \textbf{softened, solicitous, open intonation}: \textit{"Sir, aapke haath mein chot kaise lagi, zara batayein$\downarrow$"} (\texttt{[08:35]}).
\end{itemize}

\noindent\textbf{6.} \textbf{Visual \& paralinguistic performance of injury and violence: How are claims of violence substantiated non-verbally on screen? (Compare camera close-ups and pained vocal tone on police bandages/fractured hands vs. visual/acoustic treatment of student head injuries, lathi-charge footage, or defensive stammers/pauses when an officer or guest is asked who used force first.)}\par\smallskip
\begin{itemize}
  \item \textbf{Notes:}
This broadcast contains the classic textbook contrast in \textbf{multimodal substantiation of injury}:\par\smallskip
  \item \textbf{Police Injury Foregrounded (\texttt{[08:30]}--\texttt{[10:15]}):} The camera zooms into a tight close-up of the Additional DCP holding his \textbf{white-bandaged hand} up to the lens \texttt{\{\{gesture: injured DCP holds bandaged hand up to camera\}\}} while Anjana lowers her voice into a pained, sympathetic tone (\texttt{$^\circ$dekhiye DCP sahab ke haath mein fracture hai$^\circ$}, \texttt{[08:48]}) and nods gravely.
  \item \textbf{Student Injury \& Police Force Hedged (\texttt{[10:20]}--\texttt{[10:55]}):} When a reporter/question briefly asks whether police lathi-charged students first, the officer stammers and cuts himself off mid-sentence \texttt{---}: \textit{"Dekhiye wo toh--- (.) abhi investigation chal raha hai, humne sirf (.) \_preventive arrangement\_ kiya tha..."} (\texttt{[10:35]}). Instead of pressing him on the stammer \texttt{---}, Anjana nods acceptingly and pivots away \texttt{[cough / throat-clear]} (\texttt{[10:42]}), while visuals of students being dragged by police play only in a small, muted corner window.
\end{itemize}

\noindent\textbf{7.} \textbf{Asymmetric delivery of police actions vs. protester actions: Compare the vocal cadence, pauses, and hand gestures used when describing police actions (e.g., calm, matter-of-fact administrative tone on "tear gas chhode gae" / "preventive arrangement") versus protester actions (alarmist pitch, striking gestures, or ironic softening as "retaliation").}\par\smallskip
\begin{itemize}
  \item \textbf{Notes:}
  \item \textbf{Police Actions:} Described in a \textbf{calm, passive, matter-of-fact administrative cadence} with soothing open-palm gestures: \textit{"Police ki taraf se bas halka bal prayog kiya gaya (.) tear gas chhode gae taaki situation control mein rahe$\downarrow$"} (\texttt{[07:45]}), echoing the officer's euphemism \textit{"preventive arrangement"}.
  \item \textbf{Protester Actions:} Described in an \textbf{alarmist, high-velocity, active cadence} with striking hand gestures: \textit{"Unhone barricade \_tode\_! Police par \_hamla\_ kiya$\uparrow$!"} (\texttt{[07:15]}). When a student panelist claims they only pushed back in self-defence (\textit{"retaliation"}), Anjana repeats \textit{"'retaliation'$\uparrow$?!"} (\texttt{[12:25]}) with a mocking sing-song rising contour and an eye-roll \texttt{\{\{gesture: eye-roll + head toss\}\}}.
\end{itemize}

\noindent\textbf{8.} \textbf{Acoustic \& visual space granted to students' own voices: Are protesting students heard in their own voice with clear audio, steady close-ups, and empathetic listening nods, or are they reduced to noisy background chaos B-roll while studio anchors/politicians speak over them?}\par\smallskip
\begin{itemize}
  \item \textbf{Notes:}
Protesting students are \textbf{acoustically subordinated and frequently talked over \texttt{[overlap]}}. For most of the broadcast, street protesters appear in muted or noisy ambient B-roll loops showing shouting and barricade shoving while studio graphics blare \textit{"NEET KE NAAM PAR HUNGAMA"} (\texttt{[06:00]}). When a student/youth-wing representative speaks in the multi-box panel (\texttt{[11:30]}--\texttt{[13:50]}), Anjana interrupts them within 10--12 seconds \texttt{[overlap]}, raising her voice over theirs (\textit{"[Anjana: Suniye--- suniye! Barricade kisne toda?!] [Student: Ma'am humari baat toh---]"}, \texttt{[12:18]}) and eventually having the producer fade down the student's mic volume so the anchor's reprimand dominates the audio mix.\par\smallskip
\end{itemize}

\noindent\textbf{9.} \textbf{Non-verbal shift between the NEET paper-leak issue and the street protest: Does the speaker's non-verbal demeanour shift from sympathetic concern (sombre tone, sighs, solidarity gestures) when discussing the exam leak itself to combative, authoritarian, or dismissive delivery once students march toward Parliament?}\par\smallskip
\begin{itemize}
  \item \textbf{Notes:}
\textbf{Yes---a dramatic 180-degree non-verbal shift occurs.} In the opening two minutes (\texttt{[01:10]}--\texttt{[03:10]}), when mentioning the NEET paper leak in the abstract, Anjana uses a sombre, concerned tone and a sympathetic \texttt{[sigh]} (*"chhatron ke bhavishya ka sawaal hai$\downarrow$"\texttt{, }[01:45]\texttt{). However, the instant the broadcast transitions to the students marching toward Parliament (}[05:30]\texttt{ onward), her sympathy vanishes and is replaced by an \textbf{authoritarian, law-and-order prosecutorial demeanour}: her posture stiffens, her brow furrows into a scowl, her pitch spikes }$\uparrow$`, and she wags her finger at the protesters, treating collective street action as illegitimate disorder.\par\smallskip
\end{itemize}

\subsubsection*{Apply the five non-verbal features (Event 2)}

\begin{longtable}{|p{0.31\textwidth}|p{0.31\textwidth}|p{0.31\textwidth}|}
\hline
Non-verbal feature & Timestamp + spoken anchor words + non-verbal notation & What it does in context (stance, framing, implication) \\ \hline \endhead
\textbf{1. Pauses, silences \& stops} & \texttt{[10:35]} Addl. DCP: \textit{"Wo toh--- (.) abhi investigation chal raha hai"} / \texttt{[12:18]} \texttt{[overlap]} \textit{"[Anjana: Suniye--- barricade kisne toda?!] [Student: Ma'am lathi---]"} (student mic faded) & The officer's abrupt self-cut-off \texttt{---} betrays evasion when asked if police struck first; Anjana's \texttt{[overlap]} interruption and mic-fading silence the student's counter-testimony. \\ \hline
\textbf{2. Trailing sentences (\texttt{...})} & \texttt{[11:52]} \textit{"Kya ye waqai NEET ke chhatra hain (.) ya phir inke peechhe kisi siyasi party ka..."} \texttt{\{\{gesture: arched eyebrow + pointed pen\}\}} & Leaves the allegation of partisan hijacking unfinished (\texttt{...}) so the audience supplies the "Opposition conspiracy" conclusion without the anchor stating it as fact. \\ \hline
\textbf{3. Intonation, pitch \& volume} & \texttt{[08:48]} \textit{"$^\circ$Dekhiye DCP sahab ke haath mein fracture hai$^\circ$"} vs. \texttt{[06:12]} \textit{"Chhatron ke bhes mein \_UPADRAVIYON\_ ne kanoon haath mein liya$\uparrow$!"} & Softened, pained delivery \texttt{$^\circ$...$^\circ$} on police injury contrasts with loud, staccato stress and rising pitch \texttt{$\uparrow$} on \textit{"UPADRAVIYON"}, inverting victim and aggressor. \\ \hline
\textbf{4. Coughs, sighs \& vocalisations} & \texttt{[10:42]} \texttt{[[cough / throat-clear]]} (pivoting away from police lathi-charge question) / \texttt{[14:45]} \texttt{[[scoff]]} \textit{"Mudda exam hai ya siyasat$\uparrow$?!"} & The \texttt{[cough]} smooths over an awkward moment for the police guest; the \texttt{[scoff]} dismisses protesters' broader accountability demands as pure politics. \\ \hline
\textbf{5. Hand gestures, face \& posture} & \texttt{[08:40]} \texttt{\{\{gesture: injured DCP holds bandaged hand up to camera\}\}} / \texttt{[12:25]} \texttt{\{\{gesture: anchor eye-roll + finger-wag on 'retaliation'\}\}} & The bandaged hand serves as a powerful visual prop legitimising police force, while the anchor's eye-roll and finger-wag delegitimise student self-defence claims. \\ \hline
\end{longtable}

\subsubsection*{Short summary (Event 2 Non-Verbal Synthesis)}

\begin{itemize}
  \item \textbf{What overall interpretation or emotional stance does the non-verbal delivery encourage?}
The non-verbal delivery encourages \textbf{law-and-order alarm and suspicion toward street protesters}. While paying brief lip service to exam grievances, the broadcast's acoustic and visual weight guides the viewer to sympathise with the injured police officer and view the marching students as unruly puppets (\textit{"upadravi"}) of opposition parties.\par\smallskip
  \item \textbf{Whose voice or physical presence is acoustically/visually foregrounded, undermined, or silenced?}
\textbf{The injured Additional DCP (with his bandaged hand in close-up) and Anchor Anjana Om Kashyap} are visually and acoustically foregrounded, given uninterrupted time and sympathetic nods. \textbf{Protesting students} are visually reduced to chaotic B-roll and acoustically \textbf{undermined/silenced} via constant \texttt{[overlap]} interruptions, finger-wagging, and faded microphones (\texttt{[12:18]}).\par\smallskip
  \item \textbf{What evidence supports your non-verbal reading? What complicates it (e.g., speaker baseline, format)?}
The stark contrast between how Anjana treats the DCP's defensive stammer (\texttt{"wo toh--- abhi investigation..."}, \texttt{[10:35]}, met with a gentle nod and pivot) versus the student's explanation (\texttt{[12:18]}, met with an \texttt{[overlap]} shout and eye-roll) provides concrete turn-taking evidence of institutional bias. Format factors (live multi-panel TV debate time pressure) partially explain rapid interruptions, but the asymmetry in who gets interrupted proves ideological alignment.\par\smallskip
  \item \textbf{How do the non-verbal cues interact with the spoken words? (Do they reinforce, dramatise, soften/hedge, invert meaning via sarcasm/irony, or leave an accusation unspoken via trailing off?)}
The non-verbal cues \textbf{soften state violence} (calm, low-volume delivery on \textit{"tear gas chhode gae"} and \textit{"preventive arrangement"} at \texttt{[07:45]}, \texttt{[10:35]}), \textbf{dramatise protester disorder} via loud stress on \textit{"UPADRAVIYON"} (\texttt{[06:12]}) and close-ups of the DCP's bandage (\texttt{[08:40]}), \textbf{invert student claims via sarcasm} (mocking pitch on \textit{"'retaliation'$\uparrow$?!"}, \texttt{[12:25]}), and \textbf{insinuate partisan manipulation via trailing off} (\texttt{...} at \texttt{[11:52]}).\par\smallskip
  \item \textbf{What context, speaker baseline, or format factors support or complicate this interpretation? (e.g., Is this the speaker's normal style across all videos, or a marked shift on this specific actor/claim?)}
Across both Event 1 and Event 2, Anjana Om Kashyap's baseline is high-energy, authoritative Hindi TV anchoring. Yet within both videos, her delivery shifts systematically across actors: \textbf{softened, sighing, and protective toward state/temple authorities} (\texttt{$^\circ$Prabhu Shri Ram$^\circ$}, \texttt{$^\circ$DCP sahab ke haath mein fracture$^\circ$}), and \textbf{combative, high-pitched \texttt{$\uparrow$}, and accusatory toward opposition leaders and street protesters}.\par\smallskip
  \item \textbf{How does another outlet perform the same moment, actor, or claim non-verbally?}
Compare \textbf{Aaj Tak} (\texttt{Event 2}) directly with \textbf{The Wire} (\texttt{Event 2}) and \textbf{Dhruv Rathee} (\texttt{Event 2}): whereas Aaj Tak zooms in on a police officer's bandaged hand with pained vocal sympathy and cuts off protesting students \texttt{[overlap]}, The Wire and Dhruv Rathee give uninterrupted close-up microphone space to bruised and exhausted students while using skeptical rising intonation \texttt{$\uparrow$} and documentary FIRs to challenge police and NTA claims.\par\smallskip
  \item \textbf{My interpretation, with reasons:}
Aaj Tak's multimodal discourse across Events 1 and 2 operates as a \textbf{shield for institutional authority}. In Event 1, physical air-quotes (\texttt{\{\{air-quotes on 'loota'\}\}}), devotional sighs \texttt{[sigh]}, and self-repairs (\texttt{---}) protect the Ram Mandir Trust by isolating blame onto a single clerk and attacking the Opposition. In Event 2, visual props (the DCP's bandaged hand), asymmetric turn-taking \texttt{[overlap]}, and prosodic shifts from \textit{"chhatra"} to \textit{"\_UPADRAVI\_"} transform a crisis of state exam failure into a spectacle of street indiscipline.\par\smallskip
\end{itemize}

\end{document}
